% Independent mathematical audit fragment. Requires amsmath and amssymb.
% This is a relative numerical-calculus proof, not an unconditional iPRA
% calibration of the literal machine to Con(PA).
\section{AX11H: a refutation compiler for the changed halt protocol}
\label{sec:ax11h-relative}

The new command protocol combines the supplied AX11 fallthrough operation
with the supplied H test of the cached formula.  Its interaction is not
obtained by retaining the AX11 proof word unchanged: the old final Shift
is no longer a test of a visible contradiction.  This section supplies the
correct all-store compiler, treating early halts and zero aliases explicitly.

\subsection{Exact assumptions and command equations}

Let $p(a,b)=(a+b)(a+b+1)/2+b$ be Cantor pairing, and write
$[a;w]=p(a,w)$, with total projections $\operatorname{hd}$ and
$\operatorname{tl}$.  Thus
\[
 [0;0]=0,\qquad\operatorname{hd}(0)=\operatorname{tl}(0)=0,
 \qquad \operatorname{hd}([a;w])=a,
 \quad\operatorname{tl}([a;w])=w.
\]
No freeness assumption about finite lists is made.  In the inherited
coefficient-four calculus $\mathcal C_4$, implication and generalization
are coded by
\[
 I(a,b)=4p(a,b),\qquad U(i,a)=4p(a,i)+2.
\]
In particular $I(0,0)=0$, every implication code has residue zero modulo
four, and $U(i,a)\geq2$ has residue two.  Code zero has a fixed calculus
proof, and code one is the designated contradiction.

A command store is $(t,w,d,h)$: visible top, stack, offset and cached head.
The unchanged no-push modus ponens and the new Shift are
\begin{align*}
 \mathsf{MP}_x(t,w,d,h)&=
 \begin{cases}(d+x,w,d+x,0),&I(t,d+x)=h,\\
 (0,w,d+x,0),&I(t,d+x)\ne h,
 \end{cases}\\
 \mathsf S(t,w,d,h)&=
 \begin{cases}\mathsf{halt},&h=1,\\
 (h,\operatorname{tl}(w),d,\operatorname{hd}(w)),&h\ne1.
 \end{cases}
\end{align*}
An AX command has four operands.  Put $D=d+x$ and $H=h+y$.  Its exact
effect, inherited from the independently bound AX11 emitter and its full
MP-field reader, is
\[
 \mathsf A_{c,x,y,z}(t,w,d,h)
   =(0,[q_0,q_1,\ldots,q_9,t;w],z,0),                 \tag{A11H}
\]
where, in stack order,
\begin{align*}
 q_0&=G(H,D,c)\text{ or }0, & q_1&=\Gamma(c),\\
 q_2&=h_6(D,c,H),           & q_3&=h_5(D,H),\\
 q_4&=h_4(D,c,H),           & q_5&=h_3(D,c,H),\\
 q_6&=h_2(D,c,H),           & q_7&=h_1(D,H),\\
 q_8&=\rho_8(D,c,H,d),      & q_9&=U(d,h).
\end{align*}
The first choice is the guarded gate, with zero used when the gate guard
fails.  The families $G,\Gamma,h_1,\ldots,h_6,\rho_8$ are those of the
unchanged separate-diagonal calculus.  The only numerical facts about
them used by the command compiler are
\[
 q_0,\ldots,q_8\equiv0\pmod4,
 \qquad q_9=U(d,h)\equiv2\pmod4.
\]
For soundness, the required inherited interface says that each emitted
head is a theorem when $h$ is a theorem: the restricted head has its
side condition $d\leq d+x$, and the guarded gate is emitted only with
its side condition.  For completeness, the first nine entries cover
the seed families and the last entry implements theorem generalization.
This is a relative interface to $\mathcal C_4$; it is not an assumption
that $\mathcal C_4$ has already been calibrated to PA over iPRA.

The fourth operand cannot be omitted.  The fallthrough reads $z$ even
though AX has cleared the offset.  It compares $I(0,z)$ to zero.  If
$z=0$, successful MP returns zero; if $z>0$, failure returns zero.
Both outcomes retain the offset $z$ and leave the cache zero, exactly
as in (A11H).

\subsection{Stack words and the necessary early-halt case}

Write $A_0=\mathsf A_{0,0,0,0}$ and let semicolons denote execution in
left-to-right order.  Direct iteration gives
\begin{align*}
 \mathsf A_{c,x,y,0};\mathsf S^{10}(t,w,d,h)
       &=(q_8,[t;w],0,q_9),\\
 \mathsf A_{c,x,y,0};\mathsf S^{11}(t,w,d,h)
       &=(q_9,w,0,t).
\end{align*}
All eleven displayed Shifts are safe: their tested caches are zero or
one of $q_0,\ldots,q_9$, none of which is one.  The twelfth Shift tests
the old visible top $t$, which may be one.  Consequently
\[
 \mathsf{Pop}:=A_0;\mathsf S^{12}
 \quad\text{satisfies}\quad
 \mathsf{Pop}(t,w,d,h)=
 \begin{cases}
  \mathsf{halt},&t=1,\\
  (t,\operatorname{tl}(w),0,\operatorname{hd}(w)),&t\ne1.
 \end{cases}                                                   \tag{PopH}
\]
This is the essential correction to the AX11 proof.  Its Pop was
unconditional because its last sum included $q_9\geq2$.  The H cache
test has the different and useful behavior (PopH).

Define the other words by
\[
 \mathsf{Push}=A_0;\mathsf S^{10};\mathsf{MP}_0,
 \qquad \mathsf{Discard}=\mathsf{MP}_0,
 \qquad \mathsf{Offset}_i=\mathsf{MP}_i.
\]
For arbitrary natural values,
\begin{align*}
 \mathsf{Push}(t,w,d,h)&=(0,[t;w],0,0),\\
 \mathsf{Discard}(t,w,0,h)&=(0,w,0,0),\\
 \mathsf{Offset}_i(0,w,0,0)&=(0,w,i,0).
\end{align*}
Push's final comparison fails because $I(q_8,0)$ has residue zero and
$q_9$ has residue two.  Discard returns zero in either MP branch.
For Offset, the cases $i=0$ and $i>0$ join to the displayed identity.
All equations include $[0;0]=0$ and the virtual zero entries of an
empty canonical stack.

\subsection{Seed selectors and generalization}

For $0\leq j\leq8$, from any store $(t,w,0,0)$ the word
\[
 \mathsf{Discard};\mathsf A_{c,D,H,0};
       \mathsf S^{j+2};\mathsf{Pop}^{9-j}                       \tag{SelH}
\]
returns $(q_j,w,0,0)$.  After the initial Shifts the visible top is
$q_j$ and the suffix of generated heads, followed by the saved zero,
remains above the original $w$.  Every ensuing Pop is safe because
$q_j\ne1$; it removes precisely one suffix entry while preserving the
selected top.  The last Pop removes the saved zero.  Thus the whole
arbitrary old stack is protected, including all collapsed zero pairs.

To select $\rho_8(D,c,H,u)$ with $u\leq D$, use
\[
 \mathsf{Discard};\mathsf{Offset}_u;
 \mathsf A_{c,D-u,H,0};\mathsf S^{10};\mathsf{Pop}.
\]
Here $u+(D-u)=D$.  All other seed families are obtained directly from
(SelH), taking parameters satisfying their displayed guards.

The generalization word is unchanged:
\[
 \mathsf{Gen}_i=\mathsf{Push};\mathsf{Offset}_i;\mathsf S;
                  \mathsf A_{i,0,0,0};\mathsf S^{11}.
\]
It sends $(a,w,0,0)$ to $(U(i,a),w,0,0)$ for every $a,w,i$.
After Push, Offset and the single Shift the store is $(0,w,i,a)$;
that Shift tested the preceding zero cache.  AX consumes $a$ as its
generalized formula before any Shift tests that cache.  Its last emitted
head is $U(i,a)$ and its saved top is zero.  The eleven safe Shifts then
give the asserted result.  The case $a=1$ is included; no consistency
assumption or unmentioned non-one side condition is used here.

\subsection{Modus ponens and the all-store proof compiler}

Consider the suffix
\[
       \mathsf S;\mathsf S;\mathsf{MP}_b;\mathsf{Pop}            \tag{MPH}
\]
on $(0,[a,I(a,b),0;w],0,0)$.  The first Shift exposes $a$ in the
cache.  The second halts if $a=1$; otherwise it makes $a$ the top and
$I(a,b)$ the cache.  Successful MP returns
$(b,[0;w],b,0)$.  By (PopH), the final Pop halts if $b=1$ and otherwise
returns $(b,w,0,0)$.  Thus (MPH) has the exact exhaustive behavior
\[
 \begin{cases}
 \mathsf{halt},&a=1\text{ or }b=1,\\
 (b,w,0,0),&a\ne1\text{ and }b\ne1.
 \end{cases}
\]
The reserved zero is essential because successful MP itself no longer
pushes a dummy.  If $a=b=w=0$, the entire premise stack may collapse to
zero; total projections still give the nonhalting output $(0,0,0,0)$.

For a finite $\mathcal C_4$ derivation $e$, compile a word $W_e$ by
induction on its proof tree.  Use (SelH) or its restricted version for
seeds; for a generalization append $\mathsf{Gen}_i$ to the child's word.
For an MP node with major proof $e_1$ of $I(a,b)$ and minor proof $e_0$
of $a$, set
\[
 W_e=\mathsf{Discard};\mathsf{Push};W_{e_1};\mathsf{Push};
       W_{e_0};\mathsf{Push};\mathsf S;\mathsf S;
                           \mathsf{MP}_b;\mathsf{Pop}.
\]

\paragraph{Compiler invariant.}
For every valid derivation $e$, every $t$ and every $w$, execution of
$W_e$ from $(t,w,0,0)$ either halts or returns
$(\operatorname{concl}(e),w,0,0)$.

\emph{Proof.}
The selector and generalization calculations prove their cases.
In the MP case, the first Discard and Push place a reserved zero above
$w$.  The induction hypotheses protect this entire old stack when
producing the major and minor premises, and the intervening Push words
place them above that zero.  If a subword halts, the required disjunction
already holds.  Otherwise the resulting store is exactly the input of
(MPH), whose calculation gives a halt or the required conclusion while
restoring $w$.  This is a finite syntactic induction, with neither
soundness of arbitrary incoming stores nor $\operatorname{Con}(\mathrm{PA})$
as a premise.

For a derivation $e$ of one, the refuter
\[
              R_e=A_0;W_e;\mathsf{Pop}                         \tag{RefH}
\]
halts from \emph{every} store $(t,w,d,h)$.  Its initial AX resets the
offset and cache to zero and gives some finite stack.  Apply the compiler
invariant to that stack.  If no earlier halt occurred, the visible top
is one and (PopH) halts.  A final plain Shift would be incorrect: from
$(1,w,0,0)$ it does not halt and instead discards that visible one.

\subsection{Soundness, scheduling and finite witnesses}

Annotate the visible top, cache and all canonical stack entries with
$\mathcal C_4$ proofs, using the fixed proof of zero for virtual zero
entries.  Such an annotation exists for the initial zero proof store.
AX retains the old top, pushes theorem heads, and leaves theorem zero
visible and cached for every fourth operand $z$.  Successful MP combines
annotated premises; failed MP returns zero.  A nonhalting Shift moves
the corresponding annotations.  Its only halt occurs at cached code one,
whose annotation is a $\mathcal C_4$ refutation.  This proves that a
scheduler halt gives a refutation; it does not claim soundness from an
arbitrary unannotated incoming store.

The scheduler retains the four-field AX encoding.  For finite words set
$e(\epsilon)=0$ and $e(C::v)=1+b_C(e(v))$, with
\begin{align*}
 b_{\mathsf S}(z)&=p(0,z),\\
 b_{\mathsf{MP}_x}(z)&=p(1,p(x,z)),\\
 b_{\mathsf A_{c,x,y,u}}(z)&=p(c+2,p(y,p(x,p(u,z)))).
\end{align*}
The read order after the command tag is $y,x,u$.  Projection identities
recover the operands and the continuation even when zero pairs collapse.
The old three-field AX encoding would instead consume part of the next
command and is not a valid substitute.

Each positive encoded continuation is $1+b_C(z)$; its next continuation
is a component of $b_C(z)$, and hence is strictly smaller.  Therefore a
nonhalting finite candidate exhausts its command continuation.  Let
$z=e(R_e)$.  Prefixing (RefH) by $\mathsf A_{k,0,0,0}$ yields raw initial
candidate codes
\[
                 N_k=p(k+2,p(0,p(0,p(0,z)))).                  \tag{PadH}
\]
These are unbounded in $k$ and at least three.  The prefix is a total
nonhalting AX command; it may change the carried store, but the suffix
$R_e$ halts from every resulting store.  The scheduler tries consecutive
candidates beginning at two, so it reaches such a candidate unless it
has already halted.  No reset of the proof store between trials is used.

\paragraph{Primitive-recursive witness form.}
The preceding construction is stronger than an appeal to semidecision
in the standard model.  On a coded finite derivation, $e\mapsto W_e$ and
$e\mapsto N_0$ are primitive recursive.  For a proof DAG, build each line's
word from already constructed earlier words; duplicated premises enlarge
the word but do not require unbounded search.  AX is a fixed composition
of pairing and polynomial formula constructors.  MP and Shift are
primitive recursive case splits, and a candidate's number of commands
is bounded by its initial continuation plus one.  Execute those bounded
command loops for candidates $2,\ldots,N_0$, stopping on the first halt.
This is a primitive-recursive computation of a scheduler-halting witness
from a refutation witness.  Its termination follows from (RefH).

Conversely, replay a supplied finite scheduler trace and carry the above
proof annotations as additional finite data.  AX appends a fixed finite
bundle of axiom instances and, for $U(d,h)$, one generalization of the
cache's proof.  MP appends one proof step; Shift selects the corresponding
stored annotation.  The output at a halt is a derivation of one.  All
operations are primitive recursive in the finite trace.  Correctness is
proved by the compiler and annotated-store inductions just given.
Hence the command scheduler and $\mathcal C_4$ refutability have explicit
finite-witness translations in both directions.

The displayed witness maps and their correctness over the natural numbers
follow from the construction.  To obtain an iPRA proof one must put the
all-store induction into a bounded or tracked form and verify the arithmetic
coding lemmas; the modeled-scheduler frame argument provides a template.
The changing-stack universal induction has not already been reduced to
quantifier-free induction by the preceding calculations.  No new iPRA
or proof-assistant endpoint is asserted here.

\subsection{What this does and does not establish for the literal table}

To transfer the command result to the literal AX11H machine, one separately
needs its exact command/source binding, positive source-step simulation
contracts, and complete blank initialization.  Given those contracts with
primitive-recursive bounds, concatenating their traces gives a forward
finite-witness map.  For a proposed literal halt at time $T$, simulate at
most $T+1$ source primitives.  If none halted, their positive durations and
determinism would yield a literal run extending past the first halt.  This
bounded argument supplies the converse, without an unbounded inverse scan
for the next source boundary.

The new construction supplies executable bindings and literal-tape
certificates; their internalization in iPRA must still be stated and
justified at the required proof level.  Most significantly, this section
does \emph{not} establish
\[
 \mathrm{iPRA}\vdash
 \bigl(\operatorname{Refute}_{\mathcal C_4}
               \leftrightarrow\operatorname{Refute}_{\mathrm{PA}}\bigr).
\]
The current repository explicitly leaves that weak-arithmetic calibration
and parts of the physical bridge unfinished.  Composing an assumed such
calibration with the relative argument would yield the requested
$\operatorname{Con}(\mathrm{PA})$ equivalence, but an assumption must not
be relabeled as an unconditional theorem.

\subsection{A substantive protocol difference, with limited novelty}

The native table differs from AX11 in five transition rows, while retaining
194 working controls.  On the command store $(1,w,0,0)$, AX11's sum-based
Shift halts and AX11H's Shift does not.  On $(1,w,0,1)$, the former does
not halt whereas the latter does.  These are valid natural-valued command
stores, irrespective of whether they occur along a consistent blank run.
The changed behavior forces (PopH), the three MP cases, and the terminal
word (RefH); it is not control renaming, reflection or a change of initial
state.  The construction is nevertheless a new combination of existing
AX11 and H ideas, not a new principle of arithmetic calibration.

The separate exact action-graph audit compares the table against all 83
current 194-state corpus objects.  It allows every spanning root, arbitrary
state renaming, global reflection, global bit exchange, and arbitrary
terminal write/move effects.  The disconnected diagnostic has no spanning
root; all 82 complete prior graphs are compared by exact canonical tuples.
Graph novelty certifies corpus-relative structural distinction only and
does not itself prove the consistency calibration.
